We need two things from a simulator: that it produce paths of the process we specified,
and that we be able to tell whether it did.
This section gives one general algorithm, one exact algorithm for the exponential kernel,
and the test.

\subsubsection{Thinning}

\cite{LS79sim} first described an algorithm for the simulation of non-homogeneous Poisson
processes based on thinning.
\cite{Oga81lew} generalised it, replacing the requirement of bounded intensities by local
boundedness.
We recall the algorithm, since it applies to any kernel and is the control against which the
exact scheme below is checked.

The idea is Proposition \ref{prop.conditionalSurvival} run backwards.
If $\totalIntensity$ were constant, the next inter-arrival would be exponential of that rate.
It is not constant, but if it is non-increasing between events then its value at the last event
dominates it,
so one may draw an exponential at the dominating rate and keep the candidate with probability
equal to the ratio of the true rate to the dominating one.
That $\totalIntensity$ is non-increasing between events is exactly the non-negativity of the
kernel, $\excitation \geq 0$: excitation only ever decays until the next event adds to it.
The bound is therefore free, and no separate majorant has to be constructed.

We take $(\arrivalTimes[]_n,\event[n])$ as the sequence of arrival times and event labels
of the ground process, and recall that computing $\totalIntensity(t)$ requires
$\lbrace(\arrivalTimes[]_n,\event[n]): \, \arrivalTimes[]_n <t\rbrace$.

The algorithm is detailed in 
Algorithm \ref{algo.ogata}.

\begin{algorithm}[h!]
 \caption{{Ogata's thinning algorithm \citep{Oga81lew}}}
 \label{algo.ogata}
 \begin{algorithmic}[5]
  \REQUIRE $(\arrivalTimes[]_i, \event[i])_{i=1,\dots,n-1}$
  \STATE set $t:=\arrivalTimes[]_{n-1}$
  \STATE set $\xi:=0$
  \WHILE{$\xi=0$}
  \STATE draw $U \sim \text{Exp}(\totalIntensity(t))$
  \STATE set $\xi :=1$ with probability $\totalIntensity(t+U)/\totalIntensity(t)$
  \STATE update $t\leftarrow t+U$
  \ENDWHILE
  \STATE set $\arrivalTimes[]_n:=t$
  \STATE draw $\event[n]$ in $\lbrace 1,\dots,\numEventTypes\rbrace$ with probabilities proportional to $\lbrace \intensity[1](\arrivalTimes[]_n), \dots, \intensity[\numEventTypes](\arrivalTimes[]_n)\rbrace$
  \RETURN $(\arrivalTimes[]_n,\event[n])$
 \end{algorithmic}
\end{algorithm}

The state is carried from event to event by a recursion, not recomputed.
Between $\arrivalTimes[ ]_{n}$ and $\arrivalTimes[ ]_{n+1}$ the state merely decays,
so from \eqref{eq.decayedCounts},
\begin{equation}\label{eq.update_of_exponential_intensity}
 \decayedCounts(\arrivalTimes[ ]_{n+1}-)
 = \Big( \decayedCounts(\arrivalTimes[ ]_{n}-) + \unitVector_{\event[n]} \Big)
 \, e^{-\decay (\arrivalTimes[ ]_{n+1} - \arrivalTimes[ ]_{n})} .
\end{equation}
Evaluating \eqref{eq.decayedCounts} directly at each of $n$ arrival times costs $O(n^2)$,
where \eqref{eq.update_of_exponential_intensity} costs $O(1)$ per event.

\begin{remark}\label{remark.prePostJump}
 The two values in \eqref{eq.update_of_exponential_intensity} are different objects and both
 are needed.
 The pre-jump value $\decayedCounts(\arrivalTimes[ ]_{n+1}-)$ is the one that enters the
 intensity, by predictability;
 the post-jump value
 $\decayedCounts(\arrivalTimes[ ]_{n+1}-) + \unitVector_{\event[n+1]}$
 is the one carried forward as the state.
 Confusing them shifts every intensity by one event,
 and produces a path that is plausible and wrong.
\end{remark}

The algorithm is exact but wasteful: the rejected candidates are work done and discarded,
and the rejection rate grows with the branching ratio.
For the kernel of Definition \ref{def.exponentialHawkes} it can be avoided entirely.

\subsubsection{An exact scheme for the exponential kernel}

Write $\totalIntensity = \mathbf{1}^{\transpose}\intensity[ ]
= \bar{\baseIntensity} + (\mathbf{1}^{\transpose}\excitation) \decayedCounts$.
Between events every component of $\decayedCounts$ decays by the same factor,
because the decay $\decay$ is common to all pairs;
hence so does the excited part of $\totalIntensity$, and
\begin{equation}\label{eq.totalIntensityBetweenEvents}
 \totalIntensity(t) = \bar{\baseIntensity}
 + D \, e^{-\decay(t - \arrivalTimes[]_n)},
 \qquad
 D := \totalIntensity(\arrivalTimes[]_n +) - \bar{\baseIntensity},
 \qquad t \in (\arrivalTimes[]_n, \arrivalTimes[]_{n+1}].
\end{equation}
Non-negativity of $\excitation$ enters separately, and gives $D \geq 0$.

\begin{remark}\label{remark.commonDecay}
 Notice that \eqref{eq.totalIntensityBetweenEvents} uses only that $\decay$ is common to all
 pairs; constant column sums are not needed.
 With a decay $\decay\subscriptee$ depending on the pair, or even one depending only on the
 type of the exciting event, the components of $\decayedCounts$ decay at different rates,
 $\totalIntensity$ is a sum of exponentials rather than one exponential,
 and \eqref{eq.dzSurvival} below --- with it the whole scheme of this subsection --- fails.
\end{remark}

Substituting \eqref{eq.totalIntensityBetweenEvents} into
Proposition \ref{prop.conditionalSurvival} and integrating,
\begin{equation}\label{eq.dzSurvival}
 \Prob \left( \arrivalTimes[]_{n+1} - \arrivalTimes[]_{n} > s \, \vert \,
 \filtrationF_{\arrivalTimes[]_{n}} \right)
 = e^{-\bar{\baseIntensity}s} \cdot
 \exp\left( -\frac{D}{\decay}\big(1 - e^{-\decay s}\big) \right).
\end{equation}
The right-hand side is a product of two factors,
and the following lemma says what that means.

\begin{lemma}\label{lemma.minimumOfIndependent}
 For $i=1,2$ let $G_i : [0,\infty) \to [0,1]$ be non-increasing and right-continuous with
 $G_i(0)=1$, and let $\tau_i$ be independent random variables with values in
 $(0,+\infty]$ and $\Prob(\tau_i > s) = G_i(s)$.
 Then
 \begin{equation*}
  \Prob(\tau_1 \wedge \tau_2 > s) = G_1(s)G_2(s),
  \qquad
  \Prob(\tau_1 \wedge \tau_2 = +\infty) = G_1(\infty)G_2(\infty) .
 \end{equation*}
\end{lemma}
\begin{proof}
 $\lbrace \tau_1 \wedge \tau_2 > s \rbrace
 = \lbrace \tau_1 > s\rbrace \cap \lbrace \tau_2 > s\rbrace$,
 and the two events are independent.
 Letting $s \to \infty$ gives the second identity.
\end{proof}

The two factors of \eqref{eq.dzSurvival} are therefore the survival functions of two
independent waiting times, one for the immigration and one for the excitation,
and the inter-arrival time is their minimum.
Neither factor need be that of a finite random variable:
the first is $\mathrm{Exp}(\bar{\baseIntensity})$,
while the second tends to $e^{-D/\decay} > 0$, so the excitation clock is
\emph{defective} and fails to ring at all with that probability.
This is the observation of \cite{DZ13exa}.

\begin{thm}[{\citealp{DZ13exa}}]\label{thm.dassiosZhao}
 Let $\bar{\baseIntensity} > 0$, let
 $D = \totalIntensity(\arrivalTimes[]_n +) - \bar{\baseIntensity}$
 as in \eqref{eq.totalIntensityBetweenEvents},
 and let $U_1, U_2$ be independent and uniform on $(0,1)$.
 Set
 \begin{equation}\label{eq.dzDraws}
  \tau_1 := -\frac{\log U_1}{\bar{\baseIntensity}},
  \qquad
  \tau_2 := -\frac{1}{\decay}
  \log\left( 1 + \frac{\decay \log U_2}{D} \right),
 \end{equation}
 with $\tau_2 := +\infty$ when $D = 0$ or $1 + \decay \log U_2 / D \leq 0$.
 Then, conditionally on $\filtrationF_{\arrivalTimes[]_{n}}$,
 the inter-arrival time $\arrivalTimes[]_{n+1} - \arrivalTimes[]_{n}$ is distributed as
 $\tau_1 \wedge \tau_2$, and its law is \eqref{eq.dzSurvival}.
 The type is drawn at the arrival time, before the jump:
 \begin{equation}\label{eq.dzType}
  \Prob \big( \event[n+1] = e \, \vert \, \filtrationF_{\arrivalTimes[]_{n+1}-} \big)
  = \frac{\intensity(\arrivalTimes[]_{n+1}-)}{\totalIntensity(\arrivalTimes[]_{n+1}-)} .
 \end{equation}
\end{thm}
\begin{proof}
 $\tau_1$ is exponential of rate $\bar{\baseIntensity}$ by inversion.
 For $\tau_2$, write $\mathcal{X} := -\log U_2 \sim \mathrm{Exp}(1)$;
 then $\tau_2 > s$ if and only if
 $\mathcal{X} > \frac{D}{\decay}\big(1 - e^{-\decay s}\big)$,
 an event of probability
 $\exp\big(-\frac{D}{\decay}(1-e^{-\decay s})\big)$,
 which is the second factor of \eqref{eq.dzSurvival};
 and $\tau_2 = +\infty$ exactly when $\mathcal{X} \geq D/\decay$.
 The two are independent, so Lemma \ref{lemma.minimumOfIndependent} gives
 \eqref{eq.dzSurvival}.
 Equation \eqref{eq.dzType} is the definition of the intensity of the $e$-th component
 relative to that of the ground process.
\end{proof}

\begin{algorithm}[h!]
 \caption{{The exact scheme for the exponential kernel \citep{DZ13exa}}}
 \label{algo.dassiosZhao}
 \begin{algorithmic}[5]
  \REQUIRE $\arrivalTimes[]_{n}$ and the post-jump state $\decayedCounts(\arrivalTimes[]_n+)$
  \STATE set $D := \mathbf{1}^{\transpose}\excitation \decayedCounts(\arrivalTimes[]_n +)$
  \STATE draw $U_1, U_2$ independent and uniform on $(0,1)$
  \STATE set $\tau_1 := -\log U_1 / \bar{\baseIntensity}$
  \IF{$D > 0$ and $1 + \decay \log U_2 / D > 0$}
  \STATE set $\tau_2 := -\decay^{-1}\log\left(1 + \decay \log U_2 / D\right)$
  \ELSE
  \STATE set $\tau_2 := +\infty$
  \ENDIF
  \STATE set $\arrivalTimes[]_{n+1} := \arrivalTimes[]_{n} + \tau_1 \wedge \tau_2$
  \STATE set $\decayedCounts(\arrivalTimes[]_{n+1}-) :=
         \decayedCounts(\arrivalTimes[]_{n}+) \,
         e^{-\decay(\arrivalTimes[]_{n+1}-\arrivalTimes[]_{n})}$
  \STATE draw $\event[n+1]$ in $\lbrace 1,\dots,\numEventTypes\rbrace$ with probabilities
         \eqref{eq.dzType}
  \STATE set $\decayedCounts(\arrivalTimes[]_{n+1}+) :=
         \decayedCounts(\arrivalTimes[]_{n+1}-) + \unitVector_{\event[n+1]}$
  \RETURN $(\arrivalTimes[]_{n+1}, \event[n+1])$
 \end{algorithmic}
\end{algorithm}

Nothing in Algorithm \ref{algo.dassiosZhao} is approximate and nothing is rejected.

\begin{remark}\label{remark.defectiveClock}
 Notice that $\tau_2 = +\infty$ is not a numerical edge case to be patched.
 The excited factor of \eqref{eq.dzSurvival} carries the finite total mass $D/\decay$,
 which by Remark \ref{remark.countsAgainstRates} is the expected number of \emph{direct}
 offspring the current state still owes;
 with probability $e^{-D/\decay}$ it expires without firing,
 and the next event is an immigrant.
 An implementation that clips $\tau_2$ instead of returning $+\infty$ biases the immigrant
 share.
\end{remark}

\begin{remark}\label{remark.typeDraw}
 The waiting time \eqref{eq.dzDraws} sees $\excitation$ only through
 $\mathbf{1}^{\transpose}\excitation$,
 so it is \eqref{eq.dzType} that makes the scheme multivariate,
 and it is there that the full matrix re-enters.
 For the same reason $\totalIntensity$ is not itself a one-dimensional Hawkes process:
 it is exponential between events, but the size of its jump depends on the type that fires,
 so the full state $\decayedCounts$ must be carried.
\end{remark}

\subsubsection{Goodness of fit}

Theorem \ref{thm.meyer1971} turns the compensator into a test.
If the path is a path of the model, then rescaling each component's clock by its own
compensator turns it into a unit-rate Poisson process,
and the rescaled inter-arrival times \eqref{eq.time-changed_intertimes} are independent
$\text{Exp}(1)$ variables.
Both ingredients are available in closed form:
the compensator by \eqref{eq.compensatorClosedForm},
and the arrival times by construction.
The test is due to \cite{Oga88sta}, who introduced it for earthquake catalogues.

The three hypotheses are those of Theorem \ref{thm.meyer1971}:
no two components jump simultaneously;
the compensator is continuous;
and it diverges.
The second and third hold by construction for the model of this chapter.
The first holds for a simulated path, where the arrival times are almost surely distinct,
and is precisely what fails for a recorded one.

\begin{remark}\label{remark.marginsOnly}
 Comparing each component's rescaled inter-arrival times against $\text{Exp}(1)$ ---
 by a Kolmogorov--Smirnov statistic, say --- tests the \emph{marginal} law.
 Theorem \ref{thm.meyer1971} asserts two independences beyond that law:
 serial independence within a component, and joint independence across components.
 Neither is tested by a comparison of margins,
 and the second is what a violation of the no-common-jumps hypothesis breaks.
 A suite of $\numEventTypes$ passing such tests has established $\numEventTypes$ margins.
\end{remark}
